\documentclass[11pt]{article}

% Shared notes settings for Elements of AIML lectures (UPES).
% Keep this file free of \documentclass to allow reuse via \input.

\usepackage[utf8]{inputenc}
\usepackage[T1]{fontenc}
\usepackage[a4paper,margin=1in]{geometry}
\usepackage{hyperref}
\usepackage{xurl}
\usepackage{booktabs}
\usepackage{amsmath}
\usepackage{amssymb}
\usepackage{listings}
\usepackage{xcolor}
\usepackage{enumitem}
\usepackage{parskip}

% Code listings (general-purpose).
\lstset{
  basicstyle=\ttfamily\small,
  keywordstyle=\color{blue},
  commentstyle=\color{gray},
  stringstyle=\color{teal},
  showstringspaces=false,
  columns=fullflexible,
  frame=single,
  framerule=0.2pt,
  breaklines=true
}

\setlist[itemize]{topsep=0.3em,itemsep=0.2em}
\setlist[enumerate]{topsep=0.3em,itemsep=0.2em}

% Simple per-section source line.
\newcommand{\notesources}[1]{\par\noindent{\small\color{gray}\textit{Sources:} \sloppy #1\par}}

% Unit-wise source strings for this subject (used by slides/notes).
% Path is relative to the lecture `latex/` folder (the compilation CWD).
\IfFileExists{../../../_shared/aiml-sources.tex}{%
  % Source strings for Elements of AIML (used by slides + notes).

\newcommand{\AIMLRepoURL}{https://github.com/tali7c/Elements-of-AIML}

\newcommand{\AIMLLabSources}{%
Provost and Fawcett, \textit{Data Science for Business}.;%
McKinney, \textit{Python for Data Analysis}.;%
WEKA Documentation (Waikato Environment for Knowledge Analysis), accessed 2026-02-28.;%
Matplotlib and Seaborn documentation, accessed 2026-02-28.%
}

\newcommand{\AIMLUnitISources}{%
Rich and Knight, \textit{Artificial Intelligence}, McGraw Hill, 2017.;%
Russell and Norvig, \textit{Artificial Intelligence: A Modern Approach} (4e), Ch. 1--2 (Introduction, Intelligent Agents).;%
Poole and Mackworth, \textit{Artificial Intelligence: Foundations of Computational Agents} (2e), selected sections.%
}

\newcommand{\AIMLUnitIISources}{%
Russell and Norvig, \textit{Artificial Intelligence: A Modern Approach} (4e), Ch. 7--9 (Logical Agents, First-Order Logic, Inference).;%
Huth and Ryan, \textit{Logic in Computer Science} (2e), selected sections on propositional and first-order logic.;%
Genesereth and Nilsson, \textit{Logical Foundations of Artificial Intelligence}, selected chapters.%
}

\newcommand{\AIMLUnitIIISources}{%
Mueller and Massaron, \textit{Machine Learning for Dummies}, For Dummies, 2016.;%
Mitchell, \textit{Machine Learning}, selected chapters on learning basics and evaluation.;%
Scikit-learn documentation, accessed 2026-02-28.%
}

\newcommand{\AIMLUnitIVSources}{%
Mueller and Massaron, \textit{Machine Learning for Dummies}, For Dummies, 2016.;%
James, Witten, Hastie, and Tibshirani, \textit{An Introduction to Statistical Learning} (2e), selected chapters.;%
Scikit-learn documentation, accessed 2026-02-28.%
}

\newcommand{\AIMLUnitVSources}{%
NIST, \textit{AI Risk Management Framework (AI RMF 1.0)}, 2023.;%
Barocas, Hardt, and Narayanan, \textit{Fairness and Machine Learning} (online book), selected sections.;%
Knaflic, \textit{Storytelling with Data}, selected chapters on communicating results.%
}
%
}{%
  % Faculty copies live under `private/` so their compile CWD is different.
  \IfFileExists{../../../../public/_shared/aiml-sources.tex}{%
    % Source strings for Elements of AIML (used by slides + notes).

\newcommand{\AIMLRepoURL}{https://github.com/tali7c/Elements-of-AIML}

\newcommand{\AIMLLabSources}{%
Provost and Fawcett, \textit{Data Science for Business}.;%
McKinney, \textit{Python for Data Analysis}.;%
WEKA Documentation (Waikato Environment for Knowledge Analysis), accessed 2026-02-28.;%
Matplotlib and Seaborn documentation, accessed 2026-02-28.%
}

\newcommand{\AIMLUnitISources}{%
Rich and Knight, \textit{Artificial Intelligence}, McGraw Hill, 2017.;%
Russell and Norvig, \textit{Artificial Intelligence: A Modern Approach} (4e), Ch. 1--2 (Introduction, Intelligent Agents).;%
Poole and Mackworth, \textit{Artificial Intelligence: Foundations of Computational Agents} (2e), selected sections.%
}

\newcommand{\AIMLUnitIISources}{%
Russell and Norvig, \textit{Artificial Intelligence: A Modern Approach} (4e), Ch. 7--9 (Logical Agents, First-Order Logic, Inference).;%
Huth and Ryan, \textit{Logic in Computer Science} (2e), selected sections on propositional and first-order logic.;%
Genesereth and Nilsson, \textit{Logical Foundations of Artificial Intelligence}, selected chapters.%
}

\newcommand{\AIMLUnitIIISources}{%
Mueller and Massaron, \textit{Machine Learning for Dummies}, For Dummies, 2016.;%
Mitchell, \textit{Machine Learning}, selected chapters on learning basics and evaluation.;%
Scikit-learn documentation, accessed 2026-02-28.%
}

\newcommand{\AIMLUnitIVSources}{%
Mueller and Massaron, \textit{Machine Learning for Dummies}, For Dummies, 2016.;%
James, Witten, Hastie, and Tibshirani, \textit{An Introduction to Statistical Learning} (2e), selected chapters.;%
Scikit-learn documentation, accessed 2026-02-28.%
}

\newcommand{\AIMLUnitVSources}{%
NIST, \textit{AI Risk Management Framework (AI RMF 1.0)}, 2023.;%
Barocas, Hardt, and Narayanan, \textit{Fairness and Machine Learning} (online book), selected sections.;%
Knaflic, \textit{Storytelling with Data}, selected chapters on communicating results.%
}
%
  }{%
    % Source strings for Elements of AIML (used by slides + notes).

\newcommand{\AIMLRepoURL}{https://github.com/tali7c/Elements-of-AIML}

\newcommand{\AIMLLabSources}{%
Provost and Fawcett, \textit{Data Science for Business}.;%
McKinney, \textit{Python for Data Analysis}.;%
WEKA Documentation (Waikato Environment for Knowledge Analysis), accessed 2026-02-28.;%
Matplotlib and Seaborn documentation, accessed 2026-02-28.%
}

\newcommand{\AIMLUnitISources}{%
Rich and Knight, \textit{Artificial Intelligence}, McGraw Hill, 2017.;%
Russell and Norvig, \textit{Artificial Intelligence: A Modern Approach} (4e), Ch. 1--2 (Introduction, Intelligent Agents).;%
Poole and Mackworth, \textit{Artificial Intelligence: Foundations of Computational Agents} (2e), selected sections.%
}

\newcommand{\AIMLUnitIISources}{%
Russell and Norvig, \textit{Artificial Intelligence: A Modern Approach} (4e), Ch. 7--9 (Logical Agents, First-Order Logic, Inference).;%
Huth and Ryan, \textit{Logic in Computer Science} (2e), selected sections on propositional and first-order logic.;%
Genesereth and Nilsson, \textit{Logical Foundations of Artificial Intelligence}, selected chapters.%
}

\newcommand{\AIMLUnitIIISources}{%
Mueller and Massaron, \textit{Machine Learning for Dummies}, For Dummies, 2016.;%
Mitchell, \textit{Machine Learning}, selected chapters on learning basics and evaluation.;%
Scikit-learn documentation, accessed 2026-02-28.%
}

\newcommand{\AIMLUnitIVSources}{%
Mueller and Massaron, \textit{Machine Learning for Dummies}, For Dummies, 2016.;%
James, Witten, Hastie, and Tibshirani, \textit{An Introduction to Statistical Learning} (2e), selected chapters.;%
Scikit-learn documentation, accessed 2026-02-28.%
}

\newcommand{\AIMLUnitVSources}{%
NIST, \textit{AI Risk Management Framework (AI RMF 1.0)}, 2023.;%
Barocas, Hardt, and Narayanan, \textit{Fairness and Machine Learning} (online book), selected sections.;%
Knaflic, \textit{Storytelling with Data}, selected chapters on communicating results.%
}
%
  }%
}


\title{Elements of AIML\\Unit 01 -- Lecture 01 Notes\\AI Overview, Techniques \& Applications}
\author{Tofik Ali}
\date{\today}

\begin{document}
\maketitle
\tableofcontents

\section{Lecture Overview}
This lecture introduces the course and answers four practical questions:
\begin{itemize}[leftmargin=*]
  \item What do we mean by \textbf{Artificial Intelligence (AI)}?
  \item How is AI related to \textbf{Machine Learning (ML)} and \textbf{Data Science}?
  \item What are the \textbf{major technique families} used in AI systems?
  \item How do we define \textbf{success} for an AI system in the real world?
\end{itemize}
\notesources{\AIMLUnitISources}

\section{How to use these notes (self-study)}
Read the notes in order and pause at each worked example/case to ask: ``What would I do next, and why?''.
In AI, the important skill is not memorizing definitions but connecting \textbf{problem type} $\rightarrow$ \textbf{technique} $\rightarrow$ \textbf{evaluation}.
Try the practice questions before reading the answers.

\section{Learning Outcomes}
\begin{itemize}[leftmargin=*]
  \item Explain AI using multiple viewpoints (human-like vs rational, thinking vs acting).
  \item Differentiate AI vs ML vs data science using small examples.
  \item Identify symbolic, statistical, and hybrid AI techniques at a high level.
  \item Describe what a model is and list success criteria beyond accuracy.
\end{itemize}

\section{Keywords (quick definitions)}
\begin{itemize}[leftmargin=*]
  \item \textbf{Artificial Intelligence (AI):} systems that perceive, reason, learn, and act to achieve goals.
  \item \textbf{Machine Learning (ML):} methods that learn patterns from data to make predictions/decisions.
  \item \textbf{Model:} a simplified representation of reality used for prediction or reasoning.
  \item \textbf{Metric:} a quantitative measure of performance (e.g., accuracy, error, latency).
  \item \textbf{Generalization:} performance on unseen data drawn from the same process as training.
  \item \textbf{Data leakage:} unintended use of information that would not be available at prediction time.
\end{itemize}

\section{Course roadmap (how topics connect)}
This subject is intentionally organized from \textbf{foundations} to \textbf{methods} to \textbf{applications}:
\begin{itemize}[leftmargin=*]
  \item \textbf{Unit I (this unit):} what AI is, what problems it solves, and what success means in practice.
  \item \textbf{Unit II:} how to represent knowledge using logic and how to infer new facts (entailment, resolution).
  \item \textbf{Unit III:} how learning from data works (datasets, splits, cross-validation, dimensionality reduction).
  \item \textbf{Unit IV:} core ML paradigms (regression, classification, clustering, RL, semi-supervised).
  \item \textbf{Unit V:} where AI/ML is used and how to think responsibly about impacts and risks.
\end{itemize}

\section{Abbreviations and terms used in this lecture}
\begin{itemize}[leftmargin=*]
  \item \textbf{AI}: Artificial Intelligence
  \item \textbf{ML}: Machine Learning
  \item \textbf{DS}: Data Science
  \item \textbf{EDA}: Exploratory Data Analysis (lab focus)
  \item \textbf{KB}: Knowledge Base (logic-based AI; Unit II)
\end{itemize}

\section{Why AI is important now (three drivers)}
AI has existed for decades, but recent progress is strongly driven by:
\begin{itemize}[leftmargin=*]
  \item \textbf{Data availability:} many domains are now digital (transactions, text, images, sensors).
  \item \textbf{Compute availability:} GPUs/accelerators + cloud infrastructure.
  \item \textbf{Algorithm/tool improvements:} better optimization, architectures, libraries, and evaluation practices.
\end{itemize}
Practical takeaway: a strong AI system is rarely ``just the model'' --- it is usually a \textbf{full pipeline} (data + model + monitoring + feedback).

\section{AI task types (how we frame problems)}
Many AI/ML applications reduce to a small set of task types:
\begin{itemize}[leftmargin=*]
  \item \textbf{Classification:} predict a label/class (spam vs not spam).
  \item \textbf{Regression:} predict a number (house price).
  \item \textbf{Ranking/recommendation:} order items by relevance (top videos).
  \item \textbf{Clustering:} group similar items (customer segments).
  \item \textbf{Decision making:} choose actions over time (agent behavior).
\end{itemize}
Understanding the task type helps choose the right evaluation metrics and techniques.

\section{What is Artificial Intelligence?}
AI can be described using different (equally valid) perspectives. A common framing (popularized in many AI textbooks) is the 2D view:
\begin{itemize}[leftmargin=*]
  \item \textbf{Thinking vs acting} (internal reasoning vs external behavior),
  \item \textbf{Human-like vs rational} (imitate humans vs do the ``right'' thing).
\end{itemize}
This gives four informal viewpoints:
\begin{itemize}[leftmargin=*]
  \item \textbf{Thinking humanly:} model human cognition (cognitive science).
  \item \textbf{Acting humanly:} produce human-like behavior (e.g., Turing Test style).
  \item \textbf{Thinking rationally:} reason correctly using logic.
  \item \textbf{Acting rationally:} choose actions that maximize goal achievement.
\end{itemize}

\subsection{A practical course definition}
For this course, a useful working definition is:
\begin{quote}
AI is the study and engineering of systems that \textbf{perceive, reason, learn, and act} to achieve goals under uncertainty.
\end{quote}
This definition is deliberately broad: it includes both logic-based reasoning (Unit II) and data-driven learning (Units III--IV).

\subsection{Turing Test (short note)}
The \textbf{Turing Test} is a historic proposal to judge whether a machine can exhibit human-like conversational behavior.
It is influential, but it does not capture all goals of AI (for example, an AI system can be useful without being human-like).

\subsection{Rationality (preview)}
In modern AI, a central theme is \textbf{rational behavior}: choosing actions that achieve goals given limited information and resources.
We study this using the \textbf{agent} framework in Unit I Lecture 02.

\section{AI vs ML vs Data Science}
These terms are related but not identical.

\subsection{AI}
AI is the \emph{goal}: building intelligent behavior. It includes:
\begin{itemize}[leftmargin=*]
  \item \textbf{Knowledge representation and reasoning} (logic, inference, planning),
  \item \textbf{Learning} (ML),
  \item \textbf{Perception} (vision, speech),
  \item \textbf{Decision making} under uncertainty.
\end{itemize}

\subsection{Machine Learning (ML)}
ML is a \emph{toolbox} inside AI. Instead of hand-coding rules, we learn a function from data. Example tasks:
\begin{itemize}[leftmargin=*]
  \item Regression (predict a number), classification (predict a class), clustering (group data).
\end{itemize}
ML is powerful but depends heavily on data quality and correct evaluation.

\subsection{Data science}
Data science focuses on extracting insights/value from data:
\begin{itemize}[leftmargin=*]
  \item Data cleaning and exploration (EDA),
  \item Visualization and communication,
  \item Experimentation and reporting.
\end{itemize}
Often, a real project includes all three: data science to understand the problem, ML to build a predictor, and AI system design to deploy decisions safely.

\section{AI Techniques (high level)}
At a high level, techniques can be grouped into:
\begin{itemize}[leftmargin=*]
  \item \textbf{Symbolic AI:} explicit representations (facts/rules), logic, search, planning.
  \item \textbf{Statistical AI / ML:} learning from data (probabilistic and optimization-based).
  \item \textbf{Hybrid:} combinations (very common in practice).
\end{itemize}

\subsection{Symbolic AI (examples)}
Symbolic AI is useful when:
\begin{itemize}[leftmargin=*]
  \item the domain has clear rules/constraints (business policies, logic constraints),
  \item interpretability/auditability is important.
\end{itemize}
Examples of symbolic techniques:
\begin{itemize}[leftmargin=*]
  \item \textbf{Logic and inference} (Unit II): represent facts/rules and infer consequences.
  \item \textbf{Search and planning}: explore actions to reach goals (e.g., route planning).
  \item \textbf{Constraint satisfaction}: satisfy many rules simultaneously.
\end{itemize}

\subsection{Statistical AI / ML (examples)}
ML is useful when:
\begin{itemize}[leftmargin=*]
  \item rules are unknown or too complex,
  \item labeled examples are available (supervised learning), or structure exists in data (unsupervised).
\end{itemize}
Examples:
\begin{itemize}[leftmargin=*]
  \item \textbf{Regression/classification} (Unit IV): predict values/labels.
  \item \textbf{Clustering} (Unit IV): discover groups.
  \item \textbf{Dimensionality reduction} (Unit III): compress and visualize data.
\end{itemize}

\subsection{Hybrid systems}
Hybrid systems combine constraints/rules with learning. Examples:
\begin{itemize}[leftmargin=*]
  \item A fraud detection model whose decisions are filtered by business rules.
  \item A medical model that produces suggestions but must obey safety constraints and thresholds.
\end{itemize}

\subsection{When do we choose which?}
A helpful rule of thumb:
\begin{itemize}[leftmargin=*]
  \item If domain rules are known and stable, symbolic methods can be clean and interpretable.
  \item If rules are unknown or too complex, ML can learn patterns from examples.
  \item If you need both (constraints + learning), hybrid approaches are effective.
\end{itemize}

\section{Models and criteria of success}
\subsection{What is a model? (levels)}
The word ``model'' is used in multiple ways. It helps to separate levels:
\begin{itemize}[leftmargin=*]
  \item \textbf{Conceptual model:} assumptions and simplified story of the world (whiteboard).
  \item \textbf{Mathematical model:} equations/objective (e.g., loss function).
  \item \textbf{Algorithmic model:} training and inference procedure.
  \item \textbf{System model:} full pipeline (data collection, preprocessing, monitoring).
\end{itemize}

\subsection{Criteria of success}
In real deployments, accuracy alone is insufficient. Typical success criteria include:
\begin{itemize}[leftmargin=*]
  \item \textbf{Quality:} accuracy/error rates/calibration.
  \item \textbf{Cost:} compute, memory, energy, annotation effort.
  \item \textbf{Speed:} response time (latency) and throughput.
  \item \textbf{Robustness:} stability to noise and distribution shifts.
  \item \textbf{Fairness and safety:} bias, privacy, harmful outputs.
  \item \textbf{Maintainability:} monitoring, retraining, drift handling.
\end{itemize}

\subsection{Common failure modes}
\begin{itemize}[leftmargin=*]
  \item \textbf{Bad data:} missing values, mislabeled examples, sampling bias.
  \item \textbf{Leakage:} using information not available at decision time.
  \item \textbf{Train/deploy mismatch:} users and environment change.
  \item \textbf{Wrong metric:} optimizing a metric that does not reflect real cost.
  \item \textbf{No monitoring:} silent degradation due to drift.
\end{itemize}

\section{Common misconceptions (quick fixes)}
\begin{itemize}[leftmargin=*]
  \item \textbf{``AI = ML''}: ML is important, but AI also includes reasoning, planning, and agents.
  \item \textbf{``High accuracy = success''}: metric choice must match real costs and constraints.
  \item \textbf{``More data always solves it''}: data quality, bias, and leakage can still break systems.
  \item \textbf{``Model is the product''}: production systems require pipelines, monitoring, and updates.
\end{itemize}

\section{Worked example: spam filtering (rules vs learning)}
Consider the task: classify emails as \textbf{spam} or \textbf{not spam}.

\subsection{Rule-based approach (symbolic)}
We can write human-designed rules such as:
\begin{quote}
IF subject contains ``win money'' AND sender is unknown THEN spam.
\end{quote}
Strengths:
\begin{itemize}[leftmargin=*]
  \item Interpretable and easy to audit.
\end{itemize}
Weaknesses:
\begin{itemize}[leftmargin=*]
  \item Brittle: spammers adapt; rules become large and hard to maintain.
\end{itemize}

\subsection{Learning-based approach (ML)}
We collect labeled emails and learn a classifier. A typical feature set might include:
\begin{itemize}[leftmargin=*]
  \item word counts (bag-of-words), number of links, sender reputation, etc.
\end{itemize}
Strengths:
\begin{itemize}[leftmargin=*]
  \item Learns patterns that are hard to hand-code; can adapt with retraining.
\end{itemize}
Weaknesses:
\begin{itemize}[leftmargin=*]
  \item Requires data and careful evaluation; can inherit bias from labels/data.
\end{itemize}

\subsection{How do we measure success here?}
Accuracy may hide the real cost. For example:
\begin{itemize}[leftmargin=*]
  \item False positives (normal email marked as spam) can be very costly.
  \item False negatives (spam missed) are annoying but sometimes less costly.
\end{itemize}
So we might prioritize precision/recall or a cost-weighted metric (covered later in Unit IV).

\section{Common pitfalls (and how to avoid them)}
\begin{itemize}[leftmargin=*]
  \item \textbf{Treating AI, ML, and data science as the same thing:} AI is the broad goal, ML is one method family, and data science focuses on analysis/insight. Mixing them makes it harder to choose the right approach.
  \item \textbf{Thinking a model is ``just an algorithm'':} a deployed system includes data pipelines, preprocessing, monitoring, and feedback loops --- the model is only one component.
  \item \textbf{Optimizing only accuracy:} in many real tasks (fraud, medical screening, spam) different errors have different costs; choose metrics that reflect the real objective.
  \item \textbf{Ignoring generalization and leakage:} evaluation must reflect unseen data; leakage can make numbers look good and then fail in practice.
  \item \textbf{Assuming ``more complex is always better'':} simpler models can be cheaper, easier to explain, and easier to maintain (important in high-stakes domains).
\end{itemize}

\section{Practice (with answers)}

\subsection*{P1. Classify the task type.}
Decide whether each is classification, regression, clustering, or ranking:
\begin{itemize}[leftmargin=*]
  \item Predicting tomorrow's temperature.
  \item Grouping customers into segments based on purchase behavior.
  \item Ordering search results for a query.
  \item Detecting whether a transaction is fraudulent.
\end{itemize}
\textbf{Answer:} temperature = regression; customer segments = clustering; search results = ranking; fraud detection = classification (often imbalanced).

\subsection*{P2. Symbolic vs ML choice.}
Give one example where symbolic AI is a good fit and one where ML is a good fit.
\textbf{Answer:} symbolic: verifying compliance with explicit rules (policy constraints); ML: classifying spam emails from examples.

\subsection*{P3. Metric trap.}
Explain why ``accuracy'' may be a poor metric for a medical screening model.
\textbf{Answer:} because missing a true disease case (false negative) can be far more costly than a false positive; metrics like recall/sensitivity and cost-weighted measures are more appropriate.

\section{Exit questions (with answers)}

\subsection*{Q1. Give two definitions/views of AI.}
\textbf{Answer:} One view is \emph{acting rationally} (choose actions that best achieve goals). Another is \emph{thinking rationally} (reason correctly using logic). Other valid views include acting/thinking humanly.

\subsection*{Q2. Differentiate AI vs ML vs Data Science with one example each.}
\textbf{Answer:} AI is the broad goal (e.g., a rule-based expert system). ML is learning from data (e.g., a trained spam classifier). Data science is extracting insight (e.g., analyzing spam trends and presenting a report/dashboard).

\subsection*{Q3. List five criteria of success beyond accuracy.}
\textbf{Answer:} Examples: latency, compute cost/energy, robustness to shift/noise, fairness/bias, privacy/safety, maintainability/monitoring.

\subsection*{Q4. Describe one failure mode caused by data issues.}
\textbf{Answer:} Data leakage: using features that indirectly encode the label (e.g., including a post-event field that is only known after the outcome), causing unrealistically high test performance and failure in real deployment.

\section{Summary}
\begin{itemize}[leftmargin=*]
  \item AI is broad; ML is a key subset; data science supports understanding and communication.
  \item AI techniques include symbolic reasoning, ML, and hybrid approaches.
  \item Real success includes accuracy \emph{and} cost, speed, robustness, fairness, safety, and maintenance.
\end{itemize}
\notesources{\AIMLUnitISources}

\end{document}
